% TOP_PROBLEM: {"id": 20001121, "title": "Direction-complete finite-field sets and quadratic secant lifting", "category_id": 2, "category": {"id": 2, "name": "combinatorics", "display_name": "Combinatorics", "description": "Counting problems, graph theory, discrete structures.", "slug": "combinatorics", "order_index": 2, "created_at": "2026-07-31T15:26:25.671Z"}, "status": "partially_solved", "problem_number": "AIM-COMBINATORICS-0246", "set_id": 14}
% FIRST_POSTED: 2026-10-01
% ENTRY_KIND: statement_only
% UnsolvedMath ID 20001121; source catalogue code AIM-COMBINATORICS-0246
% !TeX program = lualatex
% Definition and sources only; this file is not an LLM solution attempt.
\documentclass[11pt,a4paper]{article}
\usepackage[margin=25mm]{geometry}
\usepackage{fontspec}
\setmainfont{lmroman10-regular.otf}[BoldFont=lmroman10-bold.otf,
 ItalicFont=lmroman10-italic.otf,BoldItalicFont=lmroman10-bolditalic.otf]
\usepackage{amsmath,amssymb,mathtools}
\usepackage{unicode-math}
\setmathfont{latinmodern-math.otf}
\AtBeginDocument{\let\setminus\smallsetminus}
\usepackage{xurl,hyperref}
\hypersetup{colorlinks=true,urlcolor=blue}
\setcounter{secnumdepth}{0}
\setlength{\parindent}{0pt}
\setlength{\parskip}{5pt}
\setlength{\emergencystretch}{3em}
\begin{document}
\section*{Direction-complete finite-field sets and quadratic secant lifting}\label{20001121}
{\small UnsolvedMath ID 20001121 · AIM-COMBINATORICS-0246 · Combinatorics · AIM Workshop Problem Lists}
\subsection{Definitions and mathematical statement}
Let $q$ be a prime power and let $n\geq1$.
A direction in the vector space $\mathbb F_q^n$ is a one-dimensional linear subspace
$\langle v\rangle=\{tv:t\in\mathbb F_q\}$ generated by a nonzero vector $v$.
Two nonzero vectors specify the same direction exactly when they are nonzero scalar multiples.
There are $(q^n-1)/(q-1)$ such directions.
A subset $E\subseteq\mathbb F_q^n$ determines a direction $L$ if it contains distinct points $x,y$ such that $x-y\in L\setminus\{0\}$.
The affine line through $x,y$ then has direction $L$.

Determine the extremal number
\[
 D(q,n)=\min\{|E|:E\subseteq\mathbb F_q^n,
    \text{ every direction is determined by }E\}.
\]
The field order and dimension are explicit parameters; one can seek exact values or sharp asymptotic bounds in specified regimes.
Each unordered pair of distinct points determines only one direction, so every admissible set satisfies
$\binom{|E|}{2}\geq(q^n-1)/(q-1)$.
The question is how close configurations can come to this basic counting constraint, and what additional incidence restrictions prevent equality.
The requirement is existence of a pair in every direction, not containment of a full affine line in every direction.
In particular this is a difference-direction problem rather than the finite-field Kakeya problem.
Both translations and invertible linear coordinate changes preserve admissibility and cardinality.
\subsection{Short English statement}
What is the smallest subset of a finite vector space whose point differences realize every direction?
\subsection{Sources}
\begin{itemize}
\item[S1] ULAM.AI and the UnsolvedMath Contributors, supplied problems.json, record 20001121 (AIM-COMBINATORICS-0246), AIM Workshop Problem Lists. Catalogue identity and source transcription; CC BY 4.0.
\url{https://huggingface.co/datasets/ulamai/UnsolvedMath}.
\item[S2] American Institute of Mathematics, Recent trends in additive combinatorics, item 3.9. Original workshop question underlying AIM-COMBINATORICS-0246.
\url{https://aimath.org/WWN/additivecomb/additivecomb.pdf}.
\end{itemize}
\end{document}
